\documentclass[11pt]{article}
\usepackage[margin=1in]{geometry}
\usepackage[T1]{fontenc}
\usepackage{lmodern,amsmath,amssymb,booktabs,graphicx}
\usepackage[colorlinks=true,linkcolor=blue,urlcolor=blue]{hyperref}
\setlength{\parindent}{0pt}\setlength{\parskip}{5pt}
\title{RoCE-K: joint TATE calibration and an exact-design benchmark}
\author{}\date{3 October 2026}
\begin{document}\maketitle

\textbf{Main finding.} Joint noise and nuisance control improve the existing
fitted-score procedure, but its efficiency advantage remains limited against
strong target-only inference. A separate, feasible exact-calibration procedure
in four strata gives substantially shorter protected intervals. This supports
working on score construction and empirical balance, rather than only reducing
constants in a nuisance-error certificate. It does not prove a high-dimensional
RoCE extension.

\section{Study accounting and unchanged targets}

Every site has 1,000 patients. The population target TATE is fixed at .10;
the conditional target-covariate mean is only an intermediate parameter,
and population composition uncertainty is included in the final intervals.
The DGP remains the four-stratum binary model, with target heterogeneity .3.
Weak shifts equal $2K^{-1/4}\sqrt{.5/250}$; strong shifts equal .2.
One quarter of sources depart in the treated arm. Additional scenarios put
weak deviations in different source sets across arms, or the same strong
shift in both arms so that the conditional effects cancel.
At least $.75K$ sources per arm are assumed valid throughout.

The joint-score main panel has 15 settings with 300 repetitions and an
independent confirmation panel has six settings with 1,000 repetitions.
Contrast-dispersion and exact-balance comparisons reuse those same data.
The exact-balance procedure then receives five fresh-seed settings with 500
repetitions each. A final point-rule check uses three new settings with
1,000 repetitions each. A final zero-target-heterogeneity stress panel has six settings with 500
repetitions each. There are 19,000 reported setting-specific repetitions;
paired reanalyses are not counted again. Methods within a setting share data.
All historical results and fallback cases remain available.

\section{Joint contrast certificates for the original fitted scores}

The first experiment retains the fitted weights and score center. A weighted
independent-patient empirical Bernstein bound treats the target prediction
contrast and source residual contrasts together. Source arm covariance is
computed from the same patients. A second change applies one conditional
outcome bound to the complete valid-subset nuisance contrast, instead of
adding two arm bounds.

The four combinations of separate/joint noise and separate/joint nuisance
use exactly the same error budget: .025 noise, .005 nuisance, .01 source
dispersion and .01 target reference. Source bands and reported intersections
are both saved. The target partner is fixed across the four variants.
Small-$K$ reported intervals can mostly reflect that partner, so the source
band must also be inspected.

\begin{figure}[ht]\centering
\includegraphics[width=\linewidth]{joint_ablation.pdf}
\caption{Independent weak-bias confirmation. Joint certificates shorten the
source band; the reported band also includes a target constraint. The dotted
reference uses a separately valid .95 range-adaptive target-only procedure.}
\end{figure}

At $K=512$ weak bias, the separate-noise/separate-drift source band has mean
length about .507, versus .400 with both directions joined. Reported lengths
are about .448 and .391. A strengthened range-adaptive stratified target-only
reference has length about .324. Thus the joint construction improves its
own predecessor but does not yet beat that reference. At $K=2048$ weak bias,
the joint design-certificate interval has length .337; the supplementary MGF
variant assuming source joint probabilities at least .04 gives .305.
The latter requires additional overlap information.

Direct TATE dispersion is also examined. It can use a guaranteed intersection
of arm-valid sets of size $g_\cap=g_1+g_0-K$, and it helps in the cancellation
scenario. However, this smaller valid-count guarantee increases the nuisance
certificate. In weak-bias scenarios that cost can outweigh the dispersion
improvement. The two alternatives cannot be selected by realized length unless
their joint calibration is explicitly provided.

\section{A feasible exact-calibration diagnostic}

The next comparison changes the score construction and sampling argument.
Let $\hat p_{t,x}$ be empirical proportions from all target covariates and
$N_{jxa}$ the observed source cell/arm counts. Exact design calibration uses
$q_j^a(x)=n_j\hat p_{t,x}/N_{jxa}$. For any common cell predictor,
\[
 \sum_x\hat p_{t,x}m^a(x)+\frac1{n_j}\sum_i
 \mathbf1(A_{ji}=a)q_j^a(X_{ji})\{Y_{ji}-m^a(X_{ji})\}
 =\sum_x\hat p_{t,x}\bar Y_{jxa}.
\]
The prediction cancels. Under conditional-mean transport, valid sources have
exactly the target empirical-covariate mean conditional on designs. No fitted
nuisance drift allowance or oracle variance is used.

This is a saturated standardization/calibration benchmark, not the production
sparse or fold-summed RoCE fitting program. It needs observed cell support;
a prespecified target fallback handles source cells with fewer than two
patients. Theorems and a zero-diagonal patient quadratic derivation are added
to the main design. The dispersion calibration uses unbiased cell-mean
variance corrections and Bernoulli Hoeffding proxies.

Arm dispersion remains available when arm-valid sets differ. When an
intersection is guaranteed, an adaptive option calibrates both arm and TATE
dispersion certificates simultaneously and takes the smaller protection.
The conditional source and target bands are intersected before adding the
common target population-composition allowance once.
The failure ledger is .005 for the target range, .01 for composition,
.01 each for source and target outcome noise, and .015 for dispersion.

The comparison target-only interval is strengthened again to the same
unsmoothed full-target stratified estimator, with an estimated CATE-range
certificate. It has its own .95 finite-sample guarantee. The exact-balance
procedure is therefore not compared only with the earlier, looser two-fold
AIPW reference.

\begin{center}\small
\input{balance_table}
\end{center}

These five settings each use 500 fresh repetitions. In the weak-bias setting
at $K=512$, length is .1646 with all source outcomes, versus .3083 for the
strong target-only reference. Using only the original 500 held-out source
outcomes still gives .1805. Thus extra source outcomes explain only part of
the improvement. This is not a fully isolated calibration ablation: score
construction, conditioning, and use of all target covariates also change.

The observed coverage is 1 in these settings. Finite-sample protection remains
conservative. The result establishes a useful benchmark within the saturated
model; it does not establish the same behavior for 100- or 200-dimensional
continuous covariates or for the existing RoCE estimator.

\begin{figure}[ht]\centering
\includegraphics[width=\linewidth]{conditional_balance.pdf}
\caption{Left: fresh weak-bias exact-calibration checks, with 500 repetitions
per setting. Right: a separate 1,000-repeat point-rule confirmation at $K=512$.
All methods use a fixed population truth of .10.}
\end{figure}

\section{Point estimates: a transparent tradeoff}

Protected interval coverage does not make its midpoint robust to a fixed
source departure. The point follow-up projects the target estimate onto the
conditional compatibility band. If the source conditional band contains the
finite-target conditional effect, this projection cannot increase its absolute
error relative to the target point for that conditional parameter. This is
not an unconditional population-MSE theorem.

Both point rules are evaluated on the same new-seed data, with no
scenario-specific selection of the better rule. They use exactly the same
reported confidence interval; a normal Wald interval is not attached to the
projected point.

\begin{center}\small
\input{point_table}
\end{center}

Projection substantially improves the strong-bias midpoint behavior: RMSE
changes from .0439 to .0249, versus .0312 for target-only, and bias falls
from .0417 to .0061. Under weak bias, however, its RMSE .0193 is larger than
the midpoint's .0105, while remaining below target-only. Projection is a
reasonable robust point candidate to investigate, not a uniformly optimal
choice. The implementation keeps both options explicit.

\section{Weak bias and increasing source count: a direct stress test}

The last panel removes the shared target CATE-variation floor by setting
heterogeneity to zero. The target parameter remains the population TATE .10;
this zero heterogeneity is not supplied to the inference procedure.
The invalid-source shifts still decrease as $K^{-1/4}$. All-valid controls
check an unprotected, feasible pooled normal interval using estimated
within-cell source variances and a target-composition variance estimate.
It does not correct transport bias.

\begin{center}\small
\resizebox{\linewidth}{!}{\input{weak_stress_table}}
\end{center}

The naive interval is near nominal in the all-valid controls. Its coverage
under weak bias collapses as the number of sources increases, despite smaller
individual departures. The conditional compatibility interval retains coverage
and becomes shorter. At $K=2048$ its mean length is .1020, versus .2648 for
the finite-sample-protected target-only comparator. Each row uses 500 new-seed
repetitions; overcoverage of the protected method is still visible.

The naive variance implementation is also checked against the analytic
all-valid variance in this DGP. Empirical variances from only 500 repetitions
can fluctuate; they are not substituted for the analytic or estimated variance
in any reported interval.

\begin{figure}[ht]\centering
\includegraphics[width=\linewidth]{weak_bias_stress.pdf}
\caption{Zero target effect heterogeneity exposes the weak-bias failure of
ordinary pooling. All-valid controls verify that its coverage failure is not
present when transport bias is absent. All outcomes and intervals are retained.}
\end{figure}

\section{Validation and the next theoretical task}

The checks include independent patient-level reconstructions of the weighted
mean variance, arm covariance after sample merging, a random-matrix audit of
the cell-block quadratic spectrum, exact finite-distribution calculations,
calibration cancellation, design fallback, and point-projection geometry.
Previously reported joint-score outputs are reproduced within $10^{-12}$
throughout the paired follow-up. Independent seeds are used for the new
exact-balance and point-rule checks.

The next theoretical task is to retain this conditional calibration structure
with approximate balance and high-dimensional outcome fits. Exact cancellation
of an arbitrary cellwise predictor depends on saturation. Sparse linear
outcome models offer one route through residual balancing; nonlinear links
need explicit curvature and remainder control. None of those extensions,
or a uniform MSE guarantee for the projected point, is supplied by this
finite-stratum experiment. Current RoCE/ENAR defaults and old simulation
results have not been replaced.
\end{document}
